\documentclass[11pt,a4paper]{article}

% ---------------------------------------------------------
% PREAMBLE (stable, warning-free, Overleaf-safe)
% ---------------------------------------------------------
\usepackage[utf8]{inputenc}
\usepackage[T1]{fontenc}
\usepackage{lmodern}
\usepackage{amsmath, amssymb, amsfonts}
\usepackage{bm}
\usepackage{geometry}
\usepackage{graphicx}
\usepackage{hyperref}
\usepackage{cite}
\usepackage{authblk}
\usepackage{physics}
\usepackage{mathtools}

\geometry{margin=2.4cm}

\title{\textbf{Companion XLIX: The Velocity Power Spectrum in the Relaxation-Driven Cyclic Cosmology}}
\author{Michael Lehmann \\ \texttt{mi.lehmann@gmx.de}}
\date{April 2026}

\begin{document}
\maketitle

\begin{abstract}
The velocity power spectrum provides a direct probe of the dynamical evolution of 
large-scale structure. Unlike the density power spectrum, which measures the 
amplitude of fluctuations, the velocity spectrum measures the rate at which 
structure grows and the coherence of cosmic flows. In the standard cosmological 
model, the velocity field is governed by the linear growth rate and the continuity 
equation. In the Relaxation-Driven Cyclic Cosmology (RDCC), the scale-dependent 
evolution of the gravitational potential modifies the velocity divergence, the 
amplitude of bulk flows, and the coupling between baryons and dark matter. This 
Companion develops a continuous description of the velocity power spectrum in RDCC, 
deriving how the relaxation scale influences the growth rate, the velocity 
divergence spectrum, and the observational signatures in redshift-space 
distortions. The resulting framework provides a distinctive dynamical signature of 
the RDCC gravitational field.
\end{abstract}
% ---------------------------------------------------------
\section{Introduction}

The cosmic velocity field encodes the dynamical evolution of large-scale structure. 
While the density field describes the spatial distribution of matter, the velocity 
field describes how matter flows through the Universe. These flows are driven by 
the gravitational potential and provide a direct probe of the growth rate of 
structure.

In the standard cosmological model, the velocity field is governed by the Euler and 
continuity equations, and its power spectrum is determined by the linear growth 
rate. In RDCC, the gravitational potential evolves differently. The relaxation 
mechanism suppresses the decay of small-scale modes and enhances the coherence of 
large-scale modes. This modifies the velocity divergence, the amplitude of bulk 
flows, and the coupling between baryons and dark matter. The velocity power 
spectrum becomes a direct probe of the relaxation scale $k_{\mathrm{rel}}$.
% ---------------------------------------------------------
\section{Velocity Divergence in RDCC}

The velocity divergence is defined as
\begin{equation}
    \theta(\mathbf{k},a)
    = \nabla \cdot \mathbf{v}(\mathbf{k},a).
\end{equation}

In linear theory, the continuity equation gives
\begin{equation}
    \theta(k,a)
    = -aH f(a)\,\delta(k,a),
\end{equation}
where $f(a)$ is the growth rate.

In RDCC, the density contrast evolves as
\begin{equation}
    \delta_{\mathrm{RDCC}}(k,a)
    = \delta(k,a)
      \left[
        1 - \chi_{\delta}
        \left(\frac{k}{k_{\mathrm{rel}}}\right)^{\alpha}
      \right],
\end{equation}
and the growth rate becomes scale-dependent:
\begin{equation}
    f_{\mathrm{RDCC}}(k,a)
    = f(a)
      \left[
        1 - \chi_{f}
        \left(\frac{k}{k_{\mathrm{rel}}}\right)^{\alpha}
      \right].
\end{equation}

The velocity divergence is therefore modified in a scale-dependent manner.
% ---------------------------------------------------------
\section{The RDCC Velocity Power Spectrum}

The velocity divergence power spectrum is defined as
\begin{equation}
    P_{\theta\theta}(k)
    = \langle \theta(\mathbf{k}) \theta^{*}(\mathbf{k}) \rangle.
\end{equation}

In RDCC, this becomes
\begin{equation}
    P_{\theta\theta,\mathrm{RDCC}}(k)
    = P_{\theta\theta}(k)
      \left[
        1 - \chi_{\theta}
        \left(\frac{k}{k_{\mathrm{rel}}}\right)^{\alpha}
      \right].
\end{equation}

This modification suppresses small-scale velocity power and enhances large-scale 
coherence, reflecting the relaxation mechanism.
% ---------------------------------------------------------
\section{Bulk Flows and Coherent Velocities}

Bulk flows measure the coherent motion of matter on large scales. They are given by
\begin{equation}
    \mathbf{V}(R)
    = \int d^{3}k\,
      W(kR)\,
      \mathbf{v}(\mathbf{k}),
\end{equation}
where $W$ is a window function.

In RDCC, the enhanced coherence of large-scale modes increases the amplitude of 
bulk flows:
\begin{equation}
    V_{\mathrm{RDCC}}(R)
    = V(R)
      \left[
        1 + \chi_{V}
        \left(\frac{k_{\mathrm{rel}}}{k(R)}\right)^{\alpha}
      \right].
\end{equation}

This produces a distinctive RDCC signature in peculiar velocity surveys.
% ---------------------------------------------------------
\section{Redshift-Space Distortions}

Redshift-space distortions (RSD) arise from the mapping between real-space and 
redshift-space positions. The RSD power spectrum is given by
\begin{equation}
    P_{s}(k,\mu)
    = P(k)
      \left[
        1 + f\mu^{2}
      \right]^{2},
\end{equation}
where $\mu$ is the cosine of the angle between the wave vector and the line of 
sight.

In RDCC, the growth rate becomes scale-dependent:
\begin{equation}
    f_{\mathrm{RDCC}}(k)
    = f
      \left[
        1 - \chi_{f}
        \left(\frac{k}{k_{\mathrm{rel}}}\right)^{\alpha}
      \right].
\end{equation}

The RSD signal therefore acquires a scale-dependent modification that reflects the 
relaxation scale.
% ---------------------------------------------------------
\section{Conclusion}

The velocity power spectrum in the Relaxation-Driven Cyclic Cosmology reflects the 
scale-dependent evolution of the gravitational potential. The relaxation mechanism 
modifies the velocity divergence, the amplitude of bulk flows, and the coupling 
between baryons and dark matter. These effects produce distinctive signatures in 
the velocity power spectrum, bulk flow measurements, and redshift-space 
distortions. The present Companion establishes the theoretical foundation for 
future studies of cosmic dynamics, peculiar velocities, and the growth of 
structure in RDCC.

\bibliographystyle{apsrev4-2}
\nocite{*}
\bibliography{references}

\end{document}
